\section{Recovering all part sizes}\label{sec:inverse}
A consequence of the closed formula is an inverse theorem. For example,
$K_{1,4,4}$ and $K_{2,2,5}$ have the same dimension and the same number
of edges. The coefficient $\gamma_2$ distinguishes them.

\begin{theorem}[Reconstruction]\label{thm:inverse}
Among connected complete multipartite graphs on at least two vertices,
the Ehrhart polynomial of the symmetric edge polytope determines the graph
up to isomorphism. The multiset of part sizes can be recovered by integer
arithmetic.
\end{theorem}
\begin{proof}
The degree gives $d$, hence $v=d+1$. Compute the gamma vector from the
Ehrhart series, and put $h=\lfloor d/2\rfloor$. Write
\[
 D_j=\binom d{2j}-\gamma_j/\binom{2j}{j},\qquad
 L_j(a)=\sum_{\ell=j+1}^{2j}\binom a\ell\binom{d-a}{2j-\ell}.
\]
If the cross-degree $m<h$, Theorem~\ref{thm:crossdegree} gives the
largest part $b=d-m$. It exceeds $v/2$, so it occurs once.
Record it and subtract $L_j(b)$ from each $D_j$.
All remaining parts are at most $h+1$. The same bound holds for every
part when $m=h$, since then $b\le d-h$.

For $j=h,h-1,\ldots,1$, the remaining parts now have size at most $j+1$.
The value $L_j(a)$ vanishes at every $a\le j$, while
\[
 L_j(j+1)=\binom{d-j-1}{j-1}>0 .
\]
Divide the residual $D_j$ by this number to obtain the multiplicity of
$j+1$, then subtract its contributions from the remaining equations.
The total $d+1$ determines the number of singleton parts. The algorithm
works also for $d=1$, when only this last step is needed.
\end{proof}
Since $E_G(0)=1$, the complex zero multiset with multiplicities also
determines $E_G$ and hence $G$ in this class. At order nine the two
graphs above have respectively $\gamma_2=216$ and $\gamma_2=210$.
